%%%%%%%%%%%%%%%%%%%%%%%%%%%%%%%%%%%%%%%%%%%%%%%%%%%%%%%%%%%%
%% Lecture 13
%%%%%%%%%%%%%%%%%%%%%%%%%%%%%%%%%%%%%%%%%%%%%%%%%%%%%%%%%%%%

\lecture
{13}
{Monday, 21 September 2026}
{Introduction to Probability}
{Dr. Nishant Chandgotia}


%============================================================
% Kolmogorov Extension Theorem
%============================================================

\section{Kolmogorov Extension Theorem}

\begin{theorem}{Kolmogorov's Extension Theorem}
{kolmogorov-extension-theorem}

For every $n\geq1$, let
\(
\mu_n
\)
be a probability measure on
\(
\left(\mathbb{R}^n,\mathcal{B}(\mathbb{R}^n)\right).
\)

Suppose that the family
\(
\{\mu_n:n\geq1\}
\)
is consistent in the following sense: for every $n\geq1$ and every
Borel set
\(
B\in\mathcal{B}(\mathbb{R}^n),
\)
we have
\[
\mu_{n+1}(B\times\mathbb{R})
=
\mu_n(B).
\]

Then there exists a unique probability measure
\(
\mathbb{P}
\)
on
\(
\left(
\mathbb{R}^{\mathbb{N}},
\mathcal{B}(\mathbb{R})^{\otimes\mathbb{N}}
\right)
\)
such that
\[
\mathbb{P}
\left(
B\times\mathbb{R}^{\mathbb{N}\setminus\{1,\ldots,n\}}
\right)
=
\mu_n(B)
\]
for every
\[
B\in\mathcal{B}(\mathbb{R}^n).
\]

Equivalently, for every finite-dimensional rectangle
\[
B=(a_1,b_1]\times\cdots\times(a_n,b_n],
\]
we have
\[
\mathbb{P}
\left\{
\omega\in\mathbb{R}^{\mathbb{N}}:
\omega_i\in(a_i,b_i],
\ 1\leq i\leq n
\right\}
=
\mu_n(B).
\]

\end{theorem}


\begin{remark}

The spaces $\mathbb{R}$ and $\mathbb{N}$ in the above formulation
cannot simply be replaced by arbitrary spaces without additional
assumptions.

More generally, the theorem can be formulated for suitable measurable
spaces. In particular, the relevant generalization is closely related
to standard Borel spaces.

A proof of the theorem can be found in Durrett, Theorem 2.1.21.

\end{remark}


%============================================================
% Proof
%============================================================

\subsection{Proof of the Kolmogorov Extension Theorem}

\begin{proof}

We divide the proof into several steps.

%------------------------------------------------------------
\medskip
\noindent
\textbf{Step 1: The elementary rectangles.}

Let $\mathcal{S}$ be the collection of subsets of
$\mathbb{R}^{\mathbb{N}}$ of the form
\[
A=
\left\{
\omega\in\mathbb{R}^{\mathbb{N}}:
\omega_i\in(a_i,b_i],
\quad
1\leq i\leq n
\right\},
\]
where
\[
-\infty\leq a_i<b_i\leq\infty.
\]

Equivalently,
\[
A=
(a_1,b_1]\times\cdots\times(a_n,b_n]
\times\mathbb{R}\times\mathbb{R}\times\cdots.
\]

The collection $\mathcal{S}$ is a semialgebra of subsets of
$\mathbb{R}^{\mathbb{N}}$.

Indeed, the intersection of two such sets is again of the same
form, after taking the larger of the two dimensions.

Moreover, the difference of two such rectangles can be written as a
finite disjoint union of rectangles of the same type.

%------------------------------------------------------------
\medskip
\noindent
\textbf{Step 2: Definition of $\mathbb{P}$ on $\mathcal{S}$.}

For
\[
A=
\left\{
\omega:
\omega_i\in(a_i,b_i],
\quad
1\leq i\leq n
\right\},
\]
define
\[
\boxed{
\mathbb{P}(A)
=
\mu_n
\left(
(a_1,b_1]\times\cdots\times(a_n,b_n]
\right).
}
\tag{1}
\]

We first show that this definition is well-defined.

Suppose that the same set is regarded as a subset determined by the
first $m$ coordinates, where $m>n$. Then it is represented by
\[
(a_1,b_1]\times\cdots\times(a_n,b_n]
\times\mathbb{R}^{m-n}.
\]

By consistency,
\[
\mu_m
\left(
(a_1,b_1]\times\cdots\times(a_n,b_n]
\times\mathbb{R}^{m-n}
\right)
=
\mu_n
\left(
(a_1,b_1]\times\cdots\times(a_n,b_n]
\right).
\]

Indeed, this follows by applying
\[
\mu_{k+1}(B\times\mathbb{R})
=
\mu_k(B)
\]
successively for
\[
k=n,n+1,\ldots,m-1.
\]

Thus the value in (1) does not depend on the number of coordinates
used to describe the set.

%------------------------------------------------------------
\medskip
\noindent
\textbf{Step 3: Finite additivity.}

Suppose
\[
A=A_1\sqcup\cdots\sqcup A_k,
\]
where
\[
A,A_1,\ldots,A_k\in\mathcal{S}.
\]

There exists an integer $N$ such that all these sets depend only on
the first $N$ coordinates.

Hence there are rectangles
\[
R,R_1,\ldots,R_k\subseteq\mathbb{R}^N
\]
such that
\[
A
=
R\times
\mathbb{R}^{\mathbb{N}\setminus\{1,\ldots,N\}},
\]
and similarly
\[
A_j
=
R_j\times
\mathbb{R}^{\mathbb{N}\setminus\{1,\ldots,N\}}.
\]

Moreover,
\[
R=R_1\sqcup\cdots\sqcup R_k.
\]

Since $\mu_N$ is a probability measure,
\[
\mu_N(R)
=
\sum_{j=1}^k\mu_N(R_j).
\]

Therefore,
\[
\mathbb{P}(A)
=
\sum_{j=1}^k\mathbb{P}(A_j).
\]

Thus $\mathbb{P}$ is finitely additive on $\mathcal{S}$.

%------------------------------------------------------------
\medskip
\noindent
\textbf{Step 4: Continuity lemma.}

Let $\mathcal{A}$ be the algebra generated by $\mathcal{S}$.

We claim that if
\[
B_n\in\mathcal{A},
\qquad
B_n\downarrow\varnothing,
\]
then
\[
\boxed{
\mathbb{P}(B_n)\downarrow0.
}
\tag{2}
\]

Suppose, to the contrary, that
\[
\mathbb{P}(B_n)\downarrow\delta
\]
for some
\[
\delta>0.
\]

Since each $B_n$ belongs to $\mathcal{A}$, it can be written as a
finite disjoint union of finite-dimensional rectangles. Thus, after
representing all rectangles using the first $n$ coordinates, we may
write
\[
B_n
=
\bigcup_{k=1}^{K_n}
\left\{
\omega:
\omega_i\in(a_i^k,b_i^k],
\quad
1\leq i\leq n
\right\}.
\tag{3}
\]

For every $n$, choose a finite union $C_n\subseteq B_n$ of bounded
closed rectangles such that
\[
\boxed{
\mathbb{P}(B_n\setminus C_n)
\leq
\frac{\delta}{2^{n+1}}.
}
\tag{4}
\]

Such an approximation is possible because $\mu_n$ is a probability
measure on $\mathbb{R}^n$.

Now define
\[
D_n=\bigcap_{m=1}^nC_m.
\]

Since
\[
C_m\subseteq B_m
\]
and
\[
B_n\subseteq B_m
\qquad
(m\leq n),
\]
we have
\[
D_n\subseteq B_n.
\]

Furthermore,
\[
B_n\setminus D_n
\subseteq
\bigcup_{m=1}^n(B_m\setminus C_m).
\]

Therefore, by subadditivity,
\[
\begin{aligned}
\mathbb{P}(B_n\setminus D_n)
&\leq
\sum_{m=1}^n
\mathbb{P}(B_m\setminus C_m)\\
&\leq
\sum_{m=1}^n
\frac{\delta}{2^{m+1}}\\
&<
\frac{\delta}{2}.
\end{aligned}
\]

Since
\[
\mathbb{P}(B_n)\geq\delta,
\]
we obtain
\[
\begin{aligned}
\mathbb{P}(D_n)
&=
\mathbb{P}(B_n)
-
\mathbb{P}(B_n\setminus D_n)\\
&\geq
\delta-\frac{\delta}{2}\\
&=
\frac{\delta}{2}.
\end{aligned}
\]

Thus
\[
\boxed{
\mathbb{P}(D_n)\geq\frac{\delta}{2}>0.
}
\tag{5}
\]

In particular,
\[
D_n\neq\varnothing
\]
for every $n$.

%------------------------------------------------------------
\medskip
\noindent
\textbf{Step 5: Finite-dimensional compactness argument.}

Since $D_n$ depends only on the first $n$ coordinates, there exists
a set
\[
D_n^*\subseteq\mathbb{R}^n
\]
such that
\[
D_n
=
D_n^*
\times
\mathbb{R}
\times\mathbb{R}\times\cdots.
\]

The set $D_n^*$ is a finite union of bounded closed rectangles and
hence is compact.

For each $n$, choose
\[
\omega^{(n)}\in D_n.
\]

Write
\[
\omega^{(n)}
=
\left(
\omega^{(n)}_1,
\omega^{(n)}_2,\ldots
\right).
\]

Since
\[
D_n\subseteq D_1,
\]
we have
\[
\omega^{(n)}_1\in D_1^*.
\]

By compactness of $D_1^*$, there exists a subsequence such that
\[
\omega^{(n)}_1\longrightarrow\theta_1
\]
for some
\[
\theta_1\in D_1^*.
\]

Next, since
\[
D_n\subseteq D_2
\qquad(n\geq2),
\]
we have
\[
(\omega^{(n)}_1,\omega^{(n)}_2)\in D_2^*.
\]

Using compactness of $D_2^*$, we can extract a further subsequence
such that
\[
\omega^{(n)}_2\longrightarrow\theta_2.
\]

The first coordinate still converges to $\theta_1$ because we have
only passed to a subsequence.

Continuing inductively, we obtain
\[
\theta_1,\theta_2,\ldots
\]
such that, for every $k$,
\[
\boxed{
(\theta_1,\ldots,\theta_k)\in D_k^*.
}
\tag{6}
\]

This is the diagonal-subsequence argument.

%------------------------------------------------------------
\medskip
\noindent
\textbf{Step 6: Contradiction.}

Define
\[
\theta
=
(\theta_1,\theta_2,\ldots)
\in\mathbb{R}^{\mathbb{N}}.
\]

By (6),
\[
(\theta_1,\ldots,\theta_k)\in D_k^*,
\]
and therefore
\[
\theta\in D_k
\]
for every $k$.

Consequently,
\[
\theta\in
\bigcap_{k=1}^{\infty}D_k.
\]

Since
\[
D_k\subseteq B_k,
\]
we obtain
\[
\theta\in
\bigcap_{k=1}^{\infty}B_k.
\]

Hence
\[
\bigcap_{k=1}^{\infty}B_k\neq\varnothing,
\]
which contradicts
\[
B_k\downarrow\varnothing.
\]

Therefore
\[
\mathbb{P}(B_n)\downarrow0.
\]

This proves the continuity property (2).

%------------------------------------------------------------
\medskip
\noindent
\textbf{Step 7: Countable additivity.}

Suppose
\[
A=\bigsqcup_{i=1}^{\infty}A_i,
\qquad
A,A_i\in\mathcal{S}.
\]

Define
\[
B_n
=
A\setminus\bigcup_{i=1}^nA_i.
\]

Then
\[
B_n\downarrow\varnothing.
\]

Also,
\[
A
=
A_1\sqcup\cdots\sqcup A_n\sqcup B_n.
\]

By finite additivity,
\[
\mathbb{P}(A)
=
\sum_{i=1}^n\mathbb{P}(A_i)
+
\mathbb{P}(B_n).
\]

By the continuity property,
\[
\mathbb{P}(B_n)\longrightarrow0.
\]

Therefore,
\[
\mathbb{P}(A)
=
\lim_{n\to\infty}
\sum_{i=1}^n\mathbb{P}(A_i)
=
\sum_{i=1}^{\infty}\mathbb{P}(A_i).
\]

Thus $\mathbb{P}$ is countably additive.

%------------------------------------------------------------
\medskip
\noindent
\textbf{Step 8: Extension to the product $\sigma$-field.}

We have
\[
\mathbb{P}(\mathbb{R}^{\mathbb{N}})
=
\mu_1(\mathbb{R})
=
1.
\]

Thus $\mathbb{P}$ is a finite measure on the algebra generated by
$\mathcal{S}$.

By the Carath\'eodory Extension Theorem, $\mathbb{P}$ extends to a
measure on
\[
\sigma(\mathcal{S}).
\]

But
\[
\sigma(\mathcal{S})
=
\mathcal{B}(\mathbb{R})^{\otimes\mathbb{N}}.
\]

Therefore there exists a probability measure
\[
\mathbb{P}
\]
on
\[
\left(
\mathbb{R}^{\mathbb{N}},
\mathcal{B}(\mathbb{R})^{\otimes\mathbb{N}}
\right)
\]
having the prescribed finite-dimensional distributions.

Thus existence is proved.

%------------------------------------------------------------
\medskip
\noindent
\textbf{Step 9: Uniqueness.}

Suppose that $\mathbb{P}$ and $\mathbb{Q}$ are two probability
measures on
\[
\left(
\mathbb{R}^{\mathbb{N}},
\mathcal{B}(\mathbb{R})^{\otimes\mathbb{N}}
\right)
\]
having the prescribed finite-dimensional distributions.

Then
\[
\mathbb{P}(A)=\mathbb{Q}(A)
\]
for every finite-dimensional rectangle
\[
A\in\mathcal{S}.
\]

Define
\[
\mathcal{L}
=
\left\{
A:
\mathbb{P}(A)=\mathbb{Q}(A)
\right\}.
\]

We show that $\mathcal{L}$ is a $\lambda$-system.

First,
\[
\mathbb{P}(\Omega)
=
\mathbb{Q}(\Omega)
=
1,
\]
so
\[
\Omega\in\mathcal{L}.
\]

Suppose
\[
A,B\in\mathcal{L},
\qquad
A\subseteq B.
\]

Then
\[
\mathbb{P}(B\setminus A)
=
\mathbb{P}(B)-\mathbb{P}(A)
\]
and
\[
\mathbb{Q}(B\setminus A)
=
\mathbb{Q}(B)-\mathbb{Q}(A).
\]

Since
\[
\mathbb{P}(A)=\mathbb{Q}(A)
\]
and
\[
\mathbb{P}(B)=\mathbb{Q}(B),
\]
we obtain
\[
\mathbb{P}(B\setminus A)
=
\mathbb{Q}(B\setminus A).
\]

Hence
\[
B\setminus A\in\mathcal{L}.
\]

Finally, suppose
\[
A_n\uparrow A
\]
with
\[
A_n\in\mathcal{L}.
\]

By continuity from below,
\[
\begin{aligned}
\mathbb{P}(A)
&=
\lim_{n\to\infty}\mathbb{P}(A_n)\\
&=
\lim_{n\to\infty}\mathbb{Q}(A_n)\\
&=
\mathbb{Q}(A).
\end{aligned}
\]

Thus
\[
A\in\mathcal{L}.
\]

Therefore $\mathcal{L}$ is a $\lambda$-system.

Since $\mathcal{S}$ is a $\pi$-system and
\[
\mathcal{S}\subseteq\mathcal{L},
\]
the $\pi$-$\lambda$ theorem gives
\[
\sigma(\mathcal{S})\subseteq\mathcal{L}.
\]

But
\[
\sigma(\mathcal{S})
=
\mathcal{B}(\mathbb{R})^{\otimes\mathbb{N}}.
\]

Hence
\[
\mathbb{P}(A)=\mathbb{Q}(A)
\]
for every
\[
A\in
\mathcal{B}(\mathbb{R})^{\otimes\mathbb{N}}.
\]

Therefore
\[
\mathbb{P}=\mathbb{Q}.
\]

Thus the extension is unique.

\end{proof}


%============================================================
% Remarks on the General Formulation
%============================================================

\subsection{Remarks on the General Formulation}

\begin{remark}

The space
\(
(\mathbb{R},\mathcal{B}(\mathbb{R}))
\)
may be replaced by a measurable space that is Borel isomorphic to it.

The general theory of probability measures on metric spaces provides
a broader setting for such extension results.

See, for example, the discussion of probability measures on metric
spaces in Parthasarathy.

\end{remark}


%============================================================
% Transition to Laws of Large Numbers
%============================================================

\section{Weak and Strong Laws of Large Numbers}

\subsection{Introduction}

We now move to the main limit theorems of probability theory,
beginning with the Weak Law of Large Numbers and the Strong Law of
Large Numbers.

Let
\[
X_1,X_2,\ldots
\]
be random variables and define the partial sums
\[
S_n
=
X_1+\cdots+X_n.
\]

The quantities
\[
\frac{S_n}{n}
\]
represent the empirical average of the random variables.

\begin{observation}

Under suitable assumptions, the Weak Law of Large Numbers gives
convergence of the empirical average in probability, while the Strong
Law of Large Numbers gives almost sure convergence.

In the identically distributed case, the limiting quantity is
typically
\(
\quad\mathbb{E}(X_1).
\)

\end{observation}


%============================================================
% Example: A Sequence of Random Variables
%============================================================

\subsection{Example on the Product Space}

\begin{example}{A Sequence of Random Variables on a Product Space}
{sequence-random-variables-product-space}

Consider the probability space
\[
\left(
[0,1]^{\mathbb{N}},
\mathcal{B}([0,1])^{\otimes\mathbb{N}},
\mathbb{P}
\right).
\]

For
\[
\omega=(\omega_1,\omega_2,\ldots)
\in[0,1]^{\mathbb{N}},
\]
define
\[
X_i(\omega)
=
\begin{cases}
H, & \omega_i<\omega_1,\\[4pt]
T, & \omega_i\geq\omega_1.
\end{cases}
\]

Thus each $X_i$ is a random variable defined on the product
probability space.

\end{example}


%============================================================
% Stationary Sequences and Ergodic Theorem
%============================================================

\subsection{Stationary Sequences}

\begin{definition}{Stationary Sequence of Random Variables}
{stationary-sequence}

A sequence of random variables
\(
X_1,X_2,\ldots
\)
is called \emph{stationary} if, for every
\(
n\geq1
\qquad\text{and}\qquad
k\geq0,
\)
we have
\[
(X_1,X_2,\ldots,X_n)
\overset{d}{=}
(X_{1+k},X_{2+k},\ldots,X_{n+k}).
\]

Here
\(
\overset{d}{=}
\)
means equality in distribution.

\end{definition}


\subsection{Ergodic Theorem}

\begin{observation}

Suppose that
\(
X_1,X_2,\ldots
\)
is a stationary sequence satisfying
\(
\mathbb{E}(|X_1|)<\infty.
\)

Under the assumptions of the ergodic theorem, the empirical averages
\[
\frac{S_n}{n},
\qquad where \quad S_n=X_1+\cdots+X_n,
\]
converge almost surely to a limiting random variable.

The theorem also identifies the limiting random variable in terms of
the invariant $\sigma$-field.

\end{observation}


%============================================================
% Final Note
%============================================================

\begin{profnote}

The Kolmogorov Extension Theorem is the basic tool for constructing
probability measures on infinite product spaces from consistent
finite-dimensional distributions.

The next part of the course develops the Weak and Strong Laws of
Large Numbers and then moves toward stationary processes and the
Ergodic Theorem.

\end{profnote}


%============================================================
% End of Lecture
%============================================================

\lectureend
